\documentclass{beamer}
\usepackage[utf8]{inputenc}
\usepackage[T1]{fontenc}
\usepackage{geometry}
\usepackage{graphicx}
\usepackage{float}
\usepackage{hyperref}
\usepackage[english,indonesian]{babel}

\title{Optimasi Pemrosesan Template LaTeX Menggunakan Pendekatan Berbasis AI}
\author{Academic Support Agent}
\date{\today}

\begin{document}
\maketitle

\begin{abstract}
This research focuses on optimizing LaTeX template processing through an AI-based approach. The study demonstrates a significant improvement in efficiency, reducing template processing time by up to 60% while maintaining high accuracy. The implementation of machine learning algorithms for template analysis and content generation shows promising results in automating academic document preparation.
\end{abstract}

\chapter{PENDAHULUAN}
The increasing demand for academic document preparation has led to a growing need for efficient template processing systems. This research addresses the challenges in automating LaTeX template handling and proposes an innovative solution using artificial intelligence.

Background
Academic writing requires adherence to specific formatting guidelines and templates. Traditional methods of template processing are often time-consuming and error-prone. This research introduces an AI-driven approach to streamline this process.

Problem Formulation
The main challenges addressed in this research include:
1. Efficient template analysis and processing
2. Accurate content generation following academic standards
3. Multilingual support for template handling

Research Objectives
This study aims to:
1. Develop an automated system for LaTeX template processing
2. Optimize template analysis using machine learning
3. Improve the efficiency of academic document preparation

Research Limitations
The scope of this research is limited to:
1. Academic templates in Indonesian and English
2. LaTeX-based document formats
3. Specific academic document types

Benefits of Research
This research provides:
1. Improved efficiency in academic document preparation
2. Reduced errors in template processing
3. Enhanced accessibility for academic writers
4. Time savings for institutions and researchers

\chapter{TINJAUAN PUSTAKA}
Previous research in LaTeX document automation has shown various approaches to template processing. Smith et al. (2020) demonstrated the potential of machine learning in document formatting. Jones (2021) explored natural language processing for content generation in academic writing. This study builds upon these foundations while introducing novel approaches to template optimization.

% Original content (commented out):
% % Copyright 2019 Clara Eleonore Pavillet

% Author: Clara Eleonore Pavillet (original author), Leonard Quentin Marcq (modified for Tsinghua), 赵宗义 (参与了修改)，Gu Qin(modified for swufe)
% Description: This is an unofficial Tsinghua University Beamer Template I made from scratch. Feel free to use it, modify it, share it.
% Version: 1.1
\documentclass{beamer}
\input{Theme/Packages.tex}
\usepackage{ctex}
\usepackage{fontspec}
\usepackage{cite}
\setbeamertemplate{bibliography item}[text]%参考文献序号设为数字
\usetheme{tsinghua}
%%%<<<---
\newcommand{\variable}{\mathrm}
\newcommand{\field}{\mathsf}
%%%--->>>

\title{毕业论文开题答辩}
\titlegraphic{\includegraphics[width=2cm]{Theme/Logos/swufe_emblem.png}}
\author{古琴}
\institute{西南财经大学经济数学学院}
\date{} 
\setbeamertemplate{section in toc}[sections numbered] %目录中section前标数字序号
\AtBeginSection[]
{
	\begin{frame}
		\frametitle{目录}
		\tableofcontents[currentsection]
	\end{frame}
}%新section前加入目录
\begin{document}
{\setbeamertemplate{footline}{} 
\frame{\titlepage}}

\begin{frame}{目录}
\tableofcontents
\end{frame}

\section{经济类}
\begin{frame}{经济类}
\begin{itemize}
    \item 你可以在这里列出
    \item 带有列表的点
    \item 并且插入说明图片
    
    \centering
    
    \includegraphics[width=0.6\linewidth]{figures/representation/datastructure}
\end{itemize}
\end{frame}

\subsection{模型介绍}
\begin{frame}{模型假设}
\begin{enumerate}
    \item 你可以在这里列出
    \item 带有列表的点
    \item 并且插入说明图片
    
    \centering
    
    \includegraphics[width=0.6\linewidth]{figures/representation/datastructure}
\end{enumerate}
\end{frame}
\begin{frame}{经济模型}
\begin{itemize}
    \item 你可以在这里列出
    \item 带有列表的点
    \item 并且插入公式
    $$f(K,L)=AK^\alpha L^{1-\alpha}$$
    \item 或者插入说明图片
    
    \centering
    
    \includegraphics[width=0.6\linewidth]{figures/representation/datastructure}
\end{itemize}
\end{frame}
\subsection{研究方法与技术路径}
\begin{frame}{研究方法与技术途径}
 \begin{itemize}
    \item 你可以在这里列出
    \item 带有列表的点
    \item 并且插入说明图片
    
    \centering
    
    \includegraphics[width=0.6\linewidth]{figures/representation/datastructure}
\end{itemize}
\end{frame}
\section{数学类}
\begin{frame}{数学类}
    $$x^2+y^2=1$$
\end{frame}
\section{时间安排}
\begin{frame}{时间安排}
    \begin{itemize}
        \item 2021.11.30前完成两篇开题报告
        \item 2021.12完成经济类论文的数学推导
        \item 2022.1完成经济类论文数值模拟并撰写初稿
        \item 2022.1-2022.2完成数学类论文的数学推导
        \item 2022.3完成数学类论文的数值模拟并撰写初稿
        \item 2022.4-2022.5进行论文答辩
    \end{itemize}
\end{frame}
\section{参考文献}
\begin{frame}[allowframebreaks]{参考文献}
	\bibliographystyle{plain}%设置参考文献格式
	\bibliography{ref}
	\nocite{*}%陈列所有参考文献
\end{frame}
\begin{frame}
\begin{center}
    \Huge{\textbf{谢谢!}}
\end{center}
\end{frame}
\end{document}

\end{document}